% ============================================================
%  A Theoretical Framework and Assessment Model for Evaluating
%  High-Stakes Language Test Preparation Platforms in Higher Education:
%  A Conceptual Framework Analysis
%
%  Standard LaTeX Journal Format — article class
%  Submission-Ready Version — ETHE — June 2026
% ============================================================
\documentclass[11pt,a4paper]{article}

% ── Page Layout & Margins ────────────────────────────────────
\usepackage[margin=1in]{geometry}
\usepackage{setspace}
\onehalfspacing

% ── Encoding & Fonts ─────────────────────────────────────────
\usepackage[T1]{fontenc}
\usepackage[utf8]{inputenc}
\usepackage{microtype}

% ── Mathematics ──────────────────────────────────────────────
\usepackage{amsmath,amssymb}

% ── Graphics ─────────────────────────────────────────────────
\usepackage{graphicx}
\graphicspath{{images/}}

% ── Tables ───────────────────────────────────────────────────
\usepackage{booktabs}
\usepackage{multirow}
\usepackage{array}
\usepackage{makecell}

% ── Colours ──────────────────────────────────────────────────
\usepackage{xcolor}
\definecolor{fullcol}{HTML}{2e7d32}
\definecolor{partcol}{HTML}{f57f17}
\definecolor{nonecol}{HTML}{c62828}

% ── Citations & URLs ─────────────────────────────────────────
\usepackage{cite}
\usepackage{url}

% ── Authors & Affiliations ───────────────────────────────────
\usepackage{authblk}

% ── Coverage symbols ─────────────────────────────────────────
\newcommand{\fullcov}{\textcolor{fullcol}{\textbf{\checkmark}}}
\newcommand{\partcov}{\textcolor{partcol}{\textbf{$\circ$}}}
\newcommand{\nocov}{\textcolor{nonecol}{\textbf{$\times$}}}

% ── Column types ─────────────────────────────────────────────
\newcolumntype{C}[1]{>{\centering\arraybackslash}p{#1}}
\newcolumntype{L}[1]{>{\raggedright\arraybackslash}p{#1}}

% ── Document Metadata ────────────────────────────────────────
\title{\textbf{A Theoretical Framework and Assessment Model for Evaluating High-Stakes Language Test Preparation Platforms in Higher Education: A Conceptual Framework Analysis}}

\author[1]{[Author 1 Name]\thanks{author1@institution.ac.id}}
\author[1]{[Author 2 Name]\thanks{author2@institution.ac.id}}
\affil[1]{[Department, Faculty], [Institution Name], [City, Country]}

\date{}

\begin{document}

\maketitle

% ============================================================
\begin{abstract}
The proliferation of digital technology platforms for high-stakes language test preparation (e.g., IELTS and TOEFL) in higher education has exposed a major conceptual gap. Existing evaluation frameworks are fragmented: computer-assisted language learning (CALL) studies emphasize pedagogical aspects while neglecting information systems (IS) architecture, whereas IS models prioritize technical adoption but overlook language testing specificity. To address this disciplinary divide, this study constructs an integrated theoretical framework and assessment model using Conceptual Framework Analysis (CFA) and Design Science Research (DSR) guidelines. The resulting multidimensional model integrates Technology Architecture, Pedagogy Effectiveness, and Institutional Governance dimensions, grounded in the DeLone and McLean IS Success Model as a meta-theory. The framework is operationalized into fifteen specific indicators with systematic definitions. Theoretical relationships are mapped in a structural equation model with eight propositions. The framework's utility and necessity are validated mathematically: a linear programming proof demonstrates the persistence of coverage gaps in existing frameworks, while a Monte Carlo simulation ($N = 100,000$) under Multi-Attribute Utility Theory (MAUT) complexity penalties demonstrates the Pareto dominance of the proposed framework (dominance probability of 85.80\%). A proof-of-concept inter-rater reliability test yields substantial agreement (Cohen's Kappa = 0.7917), validating the operational clarity of the assessment rubrics. This study provides educational technology researchers and university administrators with a rigorous, evidence-based instrument for evaluating and adopting high-stakes test preparation platforms.
\end{abstract}

\noindent \textbf{Keywords:} CALL evaluation, EdTech assessment framework, scoping review, dimensional gap analysis, higher education

\newpage

% ============================================================
\section{Introduction}
\label{sec:intro}
% ============================================================

\subsection{Background and Conceptual Gap}

Higher education in the digital era faces a fundamental challenge in evaluating the EdTech platforms adopted for \textit{high-stakes} language test preparation---specifically IELTS and TOEFL. The literature indicates that existing \textit{Computer-Assisted Language Learning} (CALL) platform evaluation is fragmented: it either centres on linguistic aspects without considering information systems architecture, or conversely skews too technical without interrogating pedagogical validity~\cite{b1}\cite{b2}. This gap is compounded by the absence of any evaluation framework explicitly designed for the \textit{test preparation} context, which demands precision metrics for \textit{micro-learning}, \textit{spaced repetition}, and alignment with international test standards~\cite{b3}.

A survey of the \textit{International Journal of Educational Technology in Higher Education} (ETHE) and other leading \textit{educational technology} journals reveals that conceptual research developing \textit{theoretical frameworks} and \textit{assessment models} for self-directed learning technologies receives significant attention---particularly when it bridges the disciplines of Information Systems and Applied Linguistics~\cite{b4}\cite{b5}. Yet to date, no \textit{theoretical framework} explicitly integrates the technology, pedagogy, and institutional dimensions into a single evaluation model for language test preparation platforms.

\subsection{Conceptual Research Questions}

Based on the gap identification above, three conceptual research questions are examined:
\begin{enumerate}
  \item How can a \textit{theoretical framework} be constructed that synthesises technology acceptance theory, computer-assisted language learning theory, and \textit{learning analytics} theory for evaluating high-stakes language test preparation platforms?
  \item How can an \textit{assessment model} be formulated---with indicators and operational definitions---capable of evaluating platforms from a multi-stakeholder perspective (learners, instructors, administrators) at the conceptual level?
  \item How can the theoretical relationships among evaluation dimensions be mapped in a conceptual structural model that is amenable to empirical testing in the future?
\end{enumerate}

\subsection{Research Objectives}

Table~\ref{tab:objectives} details the conceptual research objectives and corresponding outputs of this study.

\begin{table}[h]
  \caption{Research Objectives and Conceptual Outputs}
  \label{tab:objectives}
  \centering
  \renewcommand{\arraystretch}{1.3}
  \begin{tabular}{@{} L{2.2cm} L{12.0cm} @{} }
    \toprule
    \textbf{Objective} & \textbf{Conceptual Output} \\
    \midrule
    \textbf{Objective 1} & Synthesise multidisciplinary literature to identify core platform evaluation concepts \\
    \textbf{Objective 2} & Construct a multidimensional \textit{theoretical framework} through \textit{conceptual framework analysis} \\
    \textbf{Objective 3} & Formulate an \textit{assessment model} with evaluation rubrics and operational definitions \\
    \textbf{Objective 4} & Develop a conceptual structural model with theoretical propositions for future empirical testing \\
    \bottomrule
  \end{tabular}
\end{table}

\begin{figure*}[t]
  \centering
  \includegraphics[width=\textwidth]{fig5_research_program_architecture}
  \caption{Research Program Architecture illustrating the relationship between the scoping review and the conceptual framework development.}
  \label{fig:program_arch}
\end{figure*}


% ============================================================
\section{Methodology: CFA--DSR Integrated Approach}
\label{sec:method}
% ============================================================

\subsection{Philosophy and Approach}

This study adopts a hybrid methodological approach that integrates two complementary traditions in Information Systems (IS) research: \textit{Conceptual Framework Analysis} (CFA) as proposed by Jabareen (2009) for the framework construction phase~\cite{b6}, and \textit{Design Science Research} (DSR) as formalised by Hevner et al. (2004) for the artifact validation phase~\cite{b34}. This integration provides a stronger epistemological foundation than pure conceptual development: the resulting framework functions as an \textbf{IS artifact} designed to address a real-world evaluation gap.

In Gregor's (2006) taxonomy of IS theories, this study develops a \textbf{Type IV Theory (Explanatory and Predictive)}~\cite{b35}. Rather than merely classifying dimensions (Type I), it formalises structural equations (Section~\ref{sec:structural_model}) that predict how platform quality relates to outcomes. The study operates under a \textit{critical realism} paradigm at the ontological level, asserting that success dimensions of EdTech platforms are real, though our access to them is mediated by theory~\cite{b7}. Epistemologically, the study is post-positivist: knowledge is constructed through concept synthesis and validated through formal mathematical proofs and computational simulations that are reproducible and falsifiable.

Table~\ref{tab:philosophy} summarizes the philosophical and methodological layers of this study.

\begin{table}[h]
  \caption{Philosophical and Methodological Foundation}
  \label{tab:philosophy}
  \centering
  \renewcommand{\arraystretch}{1.3}
  \begin{tabular}{@{} L{3.5cm} L{4.5cm} L{6.5cm} @{} }
    \toprule
    \textbf{Layer} & \textbf{Position} & \textbf{Methodological Implications} \\
    \midrule
    \textbf{Ontology} & Critical realism & The framework is ontologically grounded, not arbitrary \\
    \textbf{Epistemology} & Post-positivism & CFA (construction) + DSR (validation) tandem approach \\
    \textbf{Methodology} & CFA + DSR & CFA for Phases 1--6; DSR for validation (Phase 7) \\
    \textbf{Literature Review} & Integrative review~\cite{b8} & Multi-disciplinary synthesis as conceptual data \\
    \textbf{Validation} & Math proofs \& simulation & Falsifiable validation without primary participant data \\
    \bottomrule
  \end{tabular}
\end{table}

\textbf{The Role of DSR in This Study:} In Hevner et al.'s (2004) terminology~\cite{b34}, the TECH+PED+INST framework serves as an \textbf{IS artifact}---a construct designed to solve a practical problem (the fragmentation of EdTech evaluation frameworks). DSR requires that the artifact be validated through evaluation activities that provide utility evidence. In this study, these activities comprise: (a) a coverage matrix cross-validation for completeness; (b) a formal linear programming proof for necessity; and (c) a Monte Carlo computational simulation for robustness. Thus, our methodological approach satisfies Hevner's three key DSR claims: utility, novelty, and rigor, supported by the Design Science Research Methodology (DSRM) guidelines of Peffers et al. (2007)~\cite{b36} for structuring the research process.

\subsection{Eight Phases of CFA in This Study}

Following Jabareen (2009)~\cite{b6}, the framework was developed through eight iterative phases. Table~\ref{tab:cfa_phases} details the procedures and conceptual outputs for each of these eight phases.

\begin{table*}[t]
  \caption{The Eight Phases of Conceptual Framework Analysis (CFA) applied in this study}
  \label{tab:cfa_phases}
  \centering
  \renewcommand{\arraystretch}{1.3}
  \begin{tabular}{@{} L{4.0cm} L{7.5cm} L{4.5cm} @{} }
    \toprule
    \textbf{Phase} & \textbf{Procedure} & \textbf{Output} \\
    \midrule
    \textbf{1. Mapping Data Sources} & Integrative literature review~\cite{b8} across Scopus, Web of Science, ERIC, and IEEE Xplore. Inclusion criteria: frameworks/models evaluating digital learning platforms in HE, CALL, MALL, or LMS (2001--2025, English). & Corpus of frameworks for conceptual analysis \\
    \textbf{2. Extensive Reading \& Categorization} & Classifying literature by discipline of origin (applied linguistics, IS, education, cognitive psychology) and type of framework (evaluative, predictive, design, adoption). & Literature categorization matrix \\
    \textbf{3. Identifying Concepts} & Extracting repeating core concepts: language learning potential, learner fit, perceived usefulness, cognitive load, learning analytics, institutional adoption. & List of candidate concepts \\
    \textbf{4. Deconstructing Concepts} & Component, history, and relationship analysis of concepts based on source theories (e.g., Sweller's cognitive load; Davis's TAM). & Componential definitions per concept \\
    \textbf{5. Integrating Concepts} & Synthesis of concepts from different disciplines into complementary evaluation dimensions. & Preliminary framework dimensions \\
    \textbf{6. Synthesis \& Resynthesis} & Construction of a network of interlinked concepts forming a plane of understanding of platform evaluation. & Theoretical framework v0.1 \\
    \textbf{7. Validation} & Three-layer DSR validation~\cite{b34}: Layer 1 (Theoretical) coverage matrix cross-validation; Layer 2 (Mathematical) linear programming proof (Theorem 1); Layer 3 (Computational) Monte Carlo sensitivity simulation (Theorem 2). & Validated framework v1.0, simulation results, formal proofs \\
    \textbf{8. Rethinking \& Revision} & Iterative revisions based on expert feedback and cross-mapping to ensure novelty and non-redundancy. & Final theoretical framework \\
    \bottomrule
  \end{tabular}
\end{table*}

\subsection{Integrative Review Protocol as CFA Data Source}

Following Torraco (2005)~\cite{b8}, an integrative review was used to collect conceptual data. Unlike systematic reviews designed to aggregate effect sizes, the integrative review enables conceptual synthesis across empirical and theoretical literature to support theory development. The review protocol was guided by the question: \textit{``What evaluation frameworks exist for digital educational platforms, and what assessment dimensions do they address or neglect in the context of higher education language learning?''} The search string combined keywords addressing educational technology evaluation, university context, and language learning.

\subsection{Mathematical Framework Specification}

To meet DSR and Type IV theory standards, we specify the formal mathematical structures of the framework:

\subsubsection{Formalization of Coverage Score}
Let $\mathcal{F} = \{F_1, \ldots, F_{12}\}$ represent the set of reviewed frameworks, and $\mathcal{D} = \{D_1, D_2, D_3, D_4, D_6\}$ represent the set of evaluation dimensions. The coding function $s: \mathcal{F} \times \mathcal{D} \to \{0, 0.5, 1.0\}$ yields the \textit{Coverage Score} (CS):
\begin{equation}
  \text{CS}(f) = \sum_{d \in \mathcal{D}} s(f,d), \quad \text{CS}(f) \in [0, 5.0]
\end{equation}
The proposed framework $F^*$ achieves $s(F^*, d) = 1.0$ for all $d$, resulting in $\text{CS}(F^*) = 5.0$. In comparison, the maximum score in the existing corpus is $\max_{f \in \mathcal{F}} \text{CS}(f) = 2.5$.

\subsubsection{Composite Platform Quality Score}
Overall platform quality is formulated as a weighted composite score:
\begin{equation}
  Q_{\text{platform}} = \alpha \cdot Q_{\text{TECH}} + \beta \cdot Q_{\text{PED}} + \gamma \cdot Q_{\text{INST}}
\end{equation}
where $\alpha + \beta + \gamma = 1$, $\alpha, \beta, \gamma > 0$, and each component score is defined as:
\begin{equation}
  Q_{\text{TECH}} = \frac{\sum_{k} w_k \cdot x_k^{\text{TECH}}}{\sum_k w_k}, \quad Q_{\text{PED}} = \frac{\sum_{j} w_j \cdot x_j^{\text{PED}}}{\sum_j w_j}, \quad Q_{\text{INST}} = \frac{\sum_{l} w_l \cdot x_l^{\text{INST}}}{\sum_l w_l}
\end{equation}
where $x_k^{\text{TECH}} \in [0,1]$ is the score of the $k$-th indicator of the TECH dimension, and similarly for the other dimensions.

\subsubsection{Formal Structural Equations Preview}
The relationships between these quality dimensions and outcomes are previewed below (detailed in Section~\ref{sec:structural_model}):
\begin{align}
  \text{PU} &= \beta_1 \cdot Q_{\text{TECH}} + \beta_2 \cdot Q_{\text{PED}} + \varepsilon_1 \\
  \text{SAT} &= \beta_3 \cdot Q_{\text{TECH}} + \beta_4 \cdot Q_{\text{PED}} + \varepsilon_2 \\
  \text{UX} &= \gamma_1 \cdot Q_{\text{TECH}} + \gamma_2 \cdot Q_{\text{PED}} + \gamma_3 \cdot (Q_{\text{TECH}} \times Q_{\text{PED}}) + \varepsilon_3 \\
  \text{USE} &= \beta_5 \cdot \text{PU} + \beta_6 \cdot \text{SAT} + \beta_7 \cdot Q_{\text{INST}} + \beta_8 \cdot \text{UX} + \varepsilon_4
\end{align}

\begin{figure*}[t]
  \centering
  \includegraphics[width=\textwidth]{fig6_pcc_and_dimensions}
  \caption{PCC Framework elements and analytical dimension grounding.}
  \label{fig:pcc_dims}
\end{figure*}


% ============================================================
\section{Literature Synthesis: Mapping Existing Frameworks}
\label{sec:literature}
% ============================================================

\subsection{Theoretical Map of CALL Evaluation}

The literature map reveals a fragmented landscape divided by disciplines. Applied Linguistics researchers focus on rich pedagogical models that ignore system architectures~\cite{b1}\cite{b2}\cite{b3}, whereas IS researchers focus on technology acceptance and technical quality while treating pedagogy as generic information quality~\cite{b4}\cite{b5}. Learning analytics frameworks address institutional monitoring but are disconnected from second language acquisition (SLA) principles~\cite{b9}\cite{b10}. 

This division creates an academic barrier: on one hand, linguists evaluate the \textit{language learning potential} of tools without assessing system architecture reliability, algorithm performance, or mobile responsiveness. On the other hand, information systems researchers employ models like TAM~\cite{b11} or UTAUT to predict adoption, yet fail to integrate variables specific to language testing. A comprehensive framework must bridge this technical-pedagogical divide.

\subsection{Conceptual Gap Analysis}

Table~\ref{tab:gaps} defines the five principal conceptual gaps identified in the existing literature.

\begin{table}[h]
  \caption{Identified Conceptual Gaps in the Literature}
  \label{tab:gaps}
  \centering
  \renewcommand{\arraystretch}{1.3}
  \begin{tabular}{@{} C{0.8cm} L{3.2cm} L{10.0cm} @{} }
    \toprule
    \textbf{No.} & \textbf{Gap} & \textbf{Manifestation in the Literature} \\
    \midrule
    1 & \textbf{Contextuality} & General CALL frameworks lack specific metrics for high-stakes test prep (test standard precision, micro-learning, spaced repetition)~\cite{b1}\cite{b3}\cite{b5}. \\
    2 & \textbf{Multidisciplinarity} & Frameworks originate from silos (linguistics vs. IS) without theoretical integration~\cite{b2}\cite{b4}. \\
    3 & \textbf{Institutionality} & Learning analytics and governance decision-support are absent from self-directed platform evaluation~\cite{b9}\cite{b10}. \\
    4 & \textbf{Structurality} & Causal path relationships between design quality, user perception, and outcomes are unmapped~\cite{b11}\cite{b12}. \\
    5 & \textbf{Operationality} & Most frameworks rely on reflective checklists rather than systematic operational definitions~\cite{b6}\cite{b13}. \\
    \bottomrule
  \end{tabular}
\end{table}

\subsubsection{Coverage Matrix: Empirical Evidence of Conceptual Gaps}

To provide empirical proof of these gaps, Table~\ref{tab:matrix} maps 12 landmark frameworks against the evaluation dimensions (TECH, PED, INST, High-Stakes specificity, Multi-Stakeholder perspective, and Rigor).

\begin{table*}[t]
  \caption{Dimensional Coverage Matrix ($\fullcov{}$ Full~=~1.0;\; $\partcov{}$ Partial~=~0.5;\; \nocov{} Absent~=~0.0.\; D5 Rigor is a separate indicator.)}
  \label{tab:matrix}
  \centering
  \renewcommand{\arraystretch}{1.3}
  \begin{tabular*}{\textwidth}{@{\extracolsep{\fill}} L{5.2cm} C{1.2cm} C{1.2cm} C{1.2cm} C{1.8cm} C{2.2cm} C{1.5cm} C{1.5cm} @{}}
    \toprule
    \textbf{Framework} & \textbf{D1 TECH} & \textbf{D2 PED} & \textbf{D3 INST} & \textbf{D4 High-Stakes} & \textbf{D6 Multi-Stakeholder} & \textbf{Score /5.0} & \textbf{D5 Rigor} \\
    \midrule
    F1 Chapelle (2001)~\cite{b1} & \nocov & \fullcov & \nocov & \partcov & \nocov & 1.5 & \partcov \\
    F2 Hubbard (2011)~\cite{b2} & \nocov & \fullcov & \nocov & \nocov & \nocov & 1.0 & \partcov \\
    F3 Leakey (2011)~\cite{b3} & \partcov & \fullcov & \nocov & \nocov & \nocov & 1.5 & \partcov \\
    F4 Colpaert (2004) & \nocov & \fullcov & \nocov & \nocov & \nocov & 1.0 & \partcov \\
    F5 Rosell-Aguilar (2017)~\cite{b5} & \partcov & \partcov & \nocov & \nocov & \nocov & 1.0 & \nocov \\
    F6 Almaiah et al. (2021) & \fullcov & \partcov & \partcov & \nocov & \partcov & 2.5 & \nocov \\
    F7 Al-Fraihat et al. (2020)~\cite{b4} & \fullcov & \partcov & \partcov & \nocov & \partcov & 2.5 & \partcov \\
    F8 Scheffel et al. (EFLA) (2014)~\cite{b9} & \nocov & \nocov & \fullcov & \nocov & \partcov & 1.5 & \partcov \\
    F9 Park \& Jo (LAD) (2019)~\cite{b10} & \nocov & \nocov & \fullcov & \nocov & \partcov & 1.5 & \partcov \\
    F10 UTAUT/Venkatesh (2003) & \partcov & \nocov & \nocov & \nocov & \nocov & 0.5 & \nocov \\
    F11 TRAM/Kampa (2024) & \partcov & \partcov & \nocov & \nocov & \nocov & 1.0 & \nocov \\
    F12 Essafi/MLLA (2025) & \partcov & \fullcov & \nocov & \partcov & \nocov & 2.0 & \nocov \\
    \midrule
    \textbf{Proposed Framework ($F^*$)} & \textbf{\fullcov} & \textbf{\fullcov} & \textbf{\fullcov} & \textbf{\fullcov} & \textbf{\fullcov} & \textbf{5.0} & \textbf{\fullcov} \\
    \bottomrule
  \end{tabular*}
\end{table*}

\subsection{Meta-Theoretical Foundation: DeLone \& McLean IS Success Model}

To avoid an arbitrary selection of dimensions, the framework is grounded in the \textbf{DeLone \& McLean IS Success Model} (D\&M, 2003)~\cite{b31}. D\&M posits that IS success depends on three quality constructs: System Quality, Information Quality, and Service Quality. These are deductively mapped to our three core dimensions as shown in Table~\ref{tab:dm_mapping}:
\begin{itemize}
  \item \textbf{Technology Architecture (TECH)} maps to \textit{System Quality} (stability, adaptivity, and interface reliability).
  \item \textbf{Pedagogy Effectiveness (PED)} maps to \textit{Information Quality} (relevance, accuracy, and utility of learning material).
  \item \textbf{Institutional Governance (INST)} maps to \textit{Service Quality} (monitoring tools, integration, and institutional support).
\end{itemize}

\begin{table}[h]
  \caption{Deductive Mapping of DeLone \& McLean (2003) Model to proposed EdTech Dimensions}
  \label{tab:dm_mapping}
  \centering
  \renewcommand{\arraystretch}{1.3}
  \begin{tabular}{@{} L{3.5cm} L{6.0cm} L{4.5cm} @{} }
    \toprule
    \textbf{D\&M Construct} & \textbf{Operationalization in Language EdTech} & \textbf{Framework Dimension} \\
    \midrule
    System Quality & Technical reliability, adaptive algorithm performance, interface layout, mobile responsiveness & Technology Architecture (TECH) \\
    Information Quality & Pedagogical content quality: standard alignment, automated diagnostics, cognitive load management, skill coverage & Pedagogy Effectiveness (PED) \\
    Service Quality & Academic governance support: progress dashboard, LMS interoperability, monitoring alerts, decision reporting & Institutional Governance (INST) \\
    Net Benefits & Individual language capability improvement (Language Learning Potential) and institutional adoption success & Outcome \\
    \bottomrule
  \end{tabular}
\end{table}

This mapping ensures that the framework's dimensions are theoretically grounded. D\&M has been empirically validated across educational technologies, showing that system, information, and service qualities jointly determine system usage and user satisfaction, which in turn drive net benefits~\cite{b32}.

\subsection{Boundary Conditions of D&M: Why CALL and SLA Theory Remain Essential}

While D\&M provides an ontological scaffold, it does not define domain-specific parameters. D\&M treats information quality generically, whereas language test preparation requires specific alignment with psychometric test format standards (e.g., IELTS rubrics). Similarly, D\&M's service quality does not capture the academic governance requirements of higher education institutions (e.g., learning analytics and Early Warning Systems). Therefore, CALL theory~\cite{b1}, SLA principles, and learning analytics frameworks~\cite{b9} are necessary to fill these domain gaps, acting as complementary extensions to D\&M. Table~\ref{tab:dm_boundaries} defines the boundary conditions and specific theoretical fillers.

\begin{table}[h]
  \caption{Boundary Conditions of DeLone \& McLean Model and Domain-Specific Extensions}
  \label{tab:dm_boundaries}
  \centering
  \renewcommand{\arraystretch}{1.3}
  \begin{tabular}{@{} L{3.5cm} L{5.0cm} L{5.5cm} @{} }
    \toprule
    \textbf{D\&M Construct} & \textbf{D\&M Boundary Limitation} & \textbf{Domain-Specific Theoretical Filler} \\
    \midrule
    Information Quality & Fails to define quality for the SLA domain & CALL theory (Chapelle, 2001)~\cite{b1}; Language Learning Potential; Washback Theory~\cite{b21} \\
    Service Quality & Does not operationalise academic governance & EFLA (Scheffel et al., 2014)~\cite{b9}; IT Governance (Tornatzky \& Fleischer, 1990)~\cite{b18} \\
    Net Benefits & Broadly defined at organizational/individual levels & Language Learning Potential + Institutional Adoption Intention \\
    \bottomrule
  \end{tabular}
\end{table}

\begin{figure*}[t]
  \centering
  \includegraphics[width=\textwidth]{fig2_coverage_matrix_heatmap}
  \caption{Heatmap of the dimensional coverage matrix showing the gaps across reviewed frameworks.}
  \label{fig:heatmap}
\end{figure*}

\begin{figure}[t]
  \centering
  \includegraphics[width=\columnwidth]{fig3_coverage_scores_ranked}
  \caption{Ranked coverage scores of reviewed frameworks.}
  \label{fig:scores}
\end{figure}

\begin{figure*}[t]
  \centering
  \includegraphics[width=\textwidth]{fig4_coverage_gaps_diagram}
  \caption{Theoretical gap diagram highlighting the unaddressed conceptual areas.}
  \label{fig:gaps_diagram}
\end{figure*}


% ============================================================
\section{Proposed Theoretical Framework}
\label{sec:framework}
% ============================================================

\subsection{Framework Philosophy}

The proposed framework is structured as a network of interlinked concepts forming a plane of understanding. The three dimensions (TECH, PED, INST) are not isolated variables but mutually articulating concepts that shape platform success. grounded in the DeLone \& McLean model (2003)~\cite{b31}, which states that system success requires system, information, and service qualities. In this domain, they are operationalised as: TECH $\leftarrow$ System Quality, PED $\leftarrow$ Information Quality, INST $\leftarrow$ Service Quality. This operationalisation is not merely a relabeling; each requires conceptual enrichment from domain-specific literature (CALL, SLA, EFLA) because D\&M does not define domain-specific metrics. This network is further enriched by crossing concepts from TAM (Davis, 1989)~\cite{b11} and Self-Determination Theory (SDT) (Deci \& Ryan, 1985)~\cite{b12} to bridge platform qualities and outcomes.

\subsection{Theoretical Justification per Dimension}

\subsubsection{Dimension 1: Technology Architecture (TECH)}
TECH captures technical system quality, drawing on Human-Computer Interaction (HCI) and Davis's TAM~\cite{b11}. System stability is grounded in software quality standards (ISO 25010)~\cite{b14}. The adaptive algorithm is grounded in intelligent tutoring systems and adaptive learning theories, stating that personalized difficulty calibration determines self-directed learning effectiveness~\cite{b15}, including prompt stability and response latency in GenAI-supported environments~\cite{b37}\cite{b40}. Minimizing extraneous cognitive load through clean interface design is grounded in cognitive load theory~\cite{b16}.

\subsubsection{Dimension 2: Pedagogy Effectiveness (PED)}
PED operationalises information quality for language learning. Materials must align with test rubrics and band descriptors (washback theory)~\cite{b21}. Micro-learning modularity matches attention constraints~\cite{b22}, while spaced repetition schedules align with the cognitive spacing effect~\cite{b23}. Automated feedback quality is grounded in Hattie \& Timperley's model, requiring timely, precise, and actionable diagnostics, especially when utilizing Large Language Models (LLMs) for speaking and writing scoring~\cite{b17}\cite{b38}\cite{b39}. Engagement mechanics are justified by Self-Determination Theory (SDT)~\cite{b12}.

\subsubsection{Dimension 3: Institutional Governance (INST)}
INST captures service quality. Dashboard analytics and monitoring tools draw on learning analytics frameworks~\cite{b9}. Interoperability requires standard API architectures~\cite{b26}. Decision-support metrics for administrators draw on IT governance and organization adoption theories~\cite{b18}. The adoption of educational technology relies on structural support, which affects adoption intention~\cite{b28}\cite{b29}.

\subsection{Cross-Cutting Concepts}

Several concepts bridge these dimensions. Table~\ref{tab:cross_cutting} defines their theoretical origins and functions.

\begin{table*}[t]
  \caption{Cross-Cutting Concepts Bridging the Core Dimensions}
  \label{tab:cross_cutting}
  \centering
  \renewcommand{\arraystretch}{1.3}
  \begin{tabular}{@{} L{4.0cm} L{4.0cm} L{8.5cm} @{} }
    \toprule
    \textbf{Concept} & \textbf{Theoretical Origin} & \textbf{Function in the Proposed Framework} \\
    \midrule
    \textbf{Language Learning Potential (LLP)} & Chapelle (2001)~\cite{b1} & Operationalises Net Benefits at the individual level. LLP serves as the core CALL evaluative criterion, mapping directly to D4 (High-Stakes Specificity). \\
    \textbf{Perceived Usefulness (PU)} & Davis (1989)~\cite{b11} & Mediates quality inputs (TECH/PED) and platform usage at the individual level. \\
    \textbf{User Experience (UX)} & Rosell-Aguilar (2017)~\cite{b5} & Captures the holistic quality of interaction between TECH and PED; mediates usage. \\
    \textbf{Practicality} & Chapelle (2001)~\cite{b1}; Hubbard (2011)~\cite{b2} & Evaluates the feasibility of the service quality (INST) for institutional adoption. \\
    \textbf{Multi-Stakeholder Perspective} & Scheffel et al. (2014)~\cite{b9} & Integrates the views of learners, instructors, and administrators; operationalised in INST-3 and INST-4. \\
    \bottomrule
  \end{tabular}
\end{table*}


% ============================================================
\section{Proposed Assessment Model}
\label{sec:assessment_model}
% ============================================================

\subsection{Assessment Model Philosophy}

\noindent The assessment model acts as a theoretical rubric with operational definitions. It is conceptual-interpretive, designed to guide evaluations and support future scale development.

\subsection{Assessment Model Hierarchy}

The proposed model establishes a three-tier hierarchy consisting of:
\begin{enumerate}
  \item \textbf{Dimensions}: The three meta-theoretical quality pillars derived from D\&M (TECH, PED, INST).
  \item \textbf{Indicators}: The 15 operationalized constructs representing specific platform capabilities (5 indicators per dimension).
  \item \textbf{Criteria}: The concrete evaluation guidelines (e.g., response latency under load, prompt stability, CEFR level precision) used by raters to assess platform performance.
\end{enumerate}

\subsection{Operational Definitions Matrix}

Table~\ref{tab:rubric} presents the indicators, operational definitions, source theories, and evaluative uses.

\begin{table*}[t]
  \caption{Assessment Model Rubric and Operational Definitions}
  \label{tab:rubric}
  \centering
  \renewcommand{\arraystretch}{1.3}
  \begin{tabular}{@{} L{2.2cm} L{5.8cm} L{3.5cm} L{3.5cm} @{} }
    \toprule
    \textbf{Indicator} & \textbf{Operational Definition} & \textbf{Source Theory} & \textbf{Evaluative Use} \\
    \midrule
    \textbf{TECH-1} Stability & Reliability of platform architecture under load & ISO 25010~\cite{b14} & Server uptime and response latency audit \\
    \textbf{TECH-2} Interface & Ease of interaction and accessibility compliance & TAM~\cite{b11}; Sweller~\cite{b16} & Cognitive walkthrough and usability evaluation \\
    \textbf{TECH-3} Adaptivity & IRT-based difficulty calibration and prompt stability & Adaptive Learning~\cite{b15}; GenAI~\cite{b37}\cite{b40} & Item difficulty alignment and prompt stability audit \\
    \textbf{TECH-4} Multimedia & Quality of audio/graphics supporting listening/reading & Dual Coding~\cite{b19}; Mayer~\cite{b20} & Content media quality assessment \\
    \textbf{TECH-5} Mobile & Responsive design and offline-first capabilities & MALL~\cite{b5} & Accessibility and offline function check \\
    \midrule
    \textbf{PED-1} Alignment & Materials aligned to IELTS/TOEFL standard rubrics & Authenticity~\cite{b1}; Washback~\cite{b21} & Item format and CEFR level matching \\
    \textbf{PED-2} Feedback & Timeliness, diagnostic precision, and AI bias detection & Hattie \& Timperley~\cite{b17}; LLMs~\cite{b38}\cite{b39} & Actionability review and AI hallucination detection \\
    \textbf{PED-3} Micro-learning & Modular units designed for 5--15 minute sessions & Sweller~\cite{b16}; Hug~\cite{b22} & Session length and structural modularity check \\
    \textbf{PED-4} Spacing & Performance-based spaced interval scheduling & Spacing Effect~\cite{b23} & Forgetting-curve algorithm logic audit \\
    \textbf{PED-5} Cognitive & Balancing extraneous and germane cognitive load & Sweller~\cite{b16} & Task interface cognitive load inspection \\
    \textbf{PED-6} Motivation & Auto-determination mechanics and anxiety reduction & SDT~\cite{b12}; Flow~\cite{b24} & Gamification and progress UI evaluation \\
    \textbf{PED-7} Skills & Commensurate coverage of reading, writing, listening, speaking & CLT~\cite{b25} & Multi-skill test matrix audit \\
    \midrule
    \textbf{INST-1} Analytics & Progress data dashboards for academic staff & EFLA~\cite{b9} & Visual representation granularity review \\
    \textbf{INST-2} Interoperability & Integration with LMS/SIAKAD via LTI/API standards & IMS Standards~\cite{b26} & Technical connection protocol audit \\
    \textbf{INST-3} Monitoring & Student tracking and Early Warning alert triggers & Early Warning~\cite{b27} & Risk alert threshold check \\
    \textbf{INST-4} Decision & Reporting cost-utility and ROI metrics for renewal & IT Governance~\cite{b18}; Diffusion~\cite{b29} & Report layout and data export evaluation \\
    \textbf{INST-5} Scalability & Capacity to scale infrastructure without latency spikes & Technology Sustainability~\cite{b28} & Load architecture review \\
    \bottomrule
  \end{tabular}
\end{table*}


% ============================================================
\section{Conceptual Structural Model}
\label{sec:structural_model}
% ============================================================

\subsection{Mapping of Theoretical Relationships}

The model maps the causal paths from inputs (TECH, PED, INST) to outcomes (Individual and Organizational Net Benefits), mediating via Perceived Usefulness, User Satisfaction, and Behavioral Intention. The moderators proficiency level, test type target, and technology readiness specify boundary effects.

Unlike early, fragmented proposals, this model positions Institutional Governance (INST) parallel to TECH and PED. Rather than predicting individual perceived usefulness, INST operates at the organizational level, determining service quality and predicting institutional adoption intention, while TECH and PED predict individual user perceptions. This separation of analytical levels prevents methodological mismatches in future empirical modeling.

\subsection{Theoretical Propositions}

Table~\ref{tab:propositions} defines the structural pathways, proposition statements, and their meta-theoretical arguments.

\begin{table*}[t]
  \caption{Structural Model Propositions and Theoretical Justifications}
  \label{tab:propositions}
  \centering
  \renewcommand{\arraystretch}{1.3}
  \begin{tabular}{@{} L{1.2cm} L{2.3cm} L{6.0cm} L{5.5cm} @{} }
    \toprule
    \textbf{Code} & \textbf{Type} & \textbf{Proposition Statement} & \textbf{Theoretical Justification} \\
    \midrule
    \textbf{P1} & Main effect & TECH quality directly increases PU and User Satisfaction. & D\&M System Quality is an antecedent of Use and Satisfaction~\cite{b31}\cite{b32}; TAM~\cite{b11}. \\
    \textbf{P2} & Main effect & PED effectiveness directly increases PU and Language Learning Potential. & D\&M Information Quality determines Satisfaction and Benefits~\cite{b31}; Chapelle's LLP~\cite{b1}. \\
    \textbf{P3} & Main effect & INST quality directly increases instructor and administrator USE, and Institutional Adoption Intention. & D\&M Service Quality operates at the organizational level, predicting institutional outcomes~\cite{b31}\cite{b32}. \\
    \textbf{P4} & Mediating pathway & PU and Satisfaction mediate the relationship between quality inputs (TECH+PED) and USE. & D\&M mediators between input quality and use behaviors~\cite{b31}; TAM~\cite{b11}. \\
    \textbf{P5} & Mediation & User Experience (UX) mediates the combined synergistic effects of TECH and PED on platform USE. & Flow Theory: optimal UX emerges from challenge-skill balance~\cite{b24}; D\&M synergy~\cite{b31}. \\
    \textbf{P6} & Net Benefits path & USE directly generates Individual Net Benefits (LLP) and student retention. & D\&M Use $\to$ Net Benefits~\cite{b31}; SDT: intrinsically motivated engagement drives competence~\cite{b12}. \\
    \textbf{P7} & Org Net Benefits & INST quality directly shapes Organizational Net Benefits through Institutional Adoption Intention. & D\&M organizational-level benefits~\cite{b31}; Diffusion of Innovations: infrastructure supports adoption~\cite{b29}. \\
    \textbf{P8} & Moderation & Initial proficiency, test type, and learner readiness moderate the structural pathways. & TTF: task-technology fit determines usage effectiveness~\cite{b30}; Washback~\cite{b21}; TRI~\cite{b33}. \\
    \bottomrule
  \end{tabular}
\end{table*}

\subsection{Formal Structural Equations}
Grounded in D\&M~\cite{b31} and SDT~\cite{b12}, the path model is formalised as:
\begin{align}
  \text{PU} &= \beta_1 \cdot Q_{\text{TECH}} + \beta_2 \cdot Q_{\text{PED}} + \varepsilon_1 \\
  \text{SAT} &= \beta_3 \cdot Q_{\text{TECH}} + \beta_4 \cdot Q_{\text{PED}} + \varepsilon_2 \\
  \text{UX} &= \gamma_1 \cdot Q_{\text{TECH}} + \gamma_2 \cdot Q_{\text{PED}} + \gamma_3 \cdot (Q_{\text{TECH}} \times Q_{\text{PED}}) + \varepsilon_3 \\
  \text{USE} &= \beta_5 \cdot \text{PU} + \beta_6 \cdot \text{SAT} + \beta_7 \cdot Q_{\text{INST}} + \beta_8 \cdot \text{UX} \nonumber \\
             &\quad + \delta_1 \cdot (\text{SAT} \times \text{PROF}) + \delta_3 \cdot (\text{PU} \times \text{TR}) + \theta_1 \cdot \text{PROF} + \theta_2 \cdot \text{TR} + \varepsilon_4 \\
  \text{NB}_{\text{ind}} &= \beta_9 \cdot \text{USE} + \delta_2 \cdot (\text{USE} \times \text{TEST}) + \theta_3 \cdot \text{TEST} + \varepsilon_5 \\
  \text{NB}_{\text{org}} &= \beta_{10} \cdot Q_{\text{INST}} + \varepsilon_6
\end{align}
where parameter $\gamma_3$ models the synergistic technology-pedagogy interaction, and $\delta_i$ model the interaction path terms.

Table~\ref{tab:parameters} presents the predicted parameter signs and their theoretical justifications.

\begin{table}[h]
  \caption{Predicted Structural Path Parameter Signs}
  \label{tab:parameters}
  \centering
  \renewcommand{\arraystretch}{1.3}
  \begin{tabular}{@{} L{3.5cm} C{2.2cm} L{7.5cm} @{} }
    \toprule
    \textbf{Parameter Path} & \textbf{Predicted Sign} & \textbf{Theoretical Basis} \\
    \midrule
    $\beta_1, \beta_2$ (PU $\leftarrow$ TECH, PED) & $> 0$ & TAM: System Quality $\to$ Perceived Usefulness~\cite{b11} \\
    $\beta_3, \beta_4$ (SAT $\leftarrow$ TECH, PED) & $> 0$ & D\&M: Quality $\to$ User Satisfaction~\cite{b31} \\
    $\gamma_1, \gamma_2, \gamma_3$ (UX $\leftarrow$ TECH, PED) & $> 0$ & Flow Theory~\cite{b24}: $\gamma_3$ represents synergy \\
    $\beta_5, \beta_6, \beta_7, \beta_8$ (USE) & $> 0$ & D\&M: PU, SAT, Service Quality $\to$ Use~\cite{b31} \\
    $\beta_9$ ($\text{NB}_{\text{ind}} \leftarrow$ USE) & $> 0$ & D\&M: Use $\to$ Net Benefits~\cite{b31}; SDT~\cite{b12} \\
    $\beta_{10}$ ($\text{NB}_{\text{org}} \leftarrow$ INST) & $> 0$ & D\&M: Service Quality $\to$ Organizational Net Benefits~\cite{b31} \\
    $\delta_1$ (SAT $\times$ PROF moderation) & $+$ & TTF: proficiency fit strengthens satisfaction-use path~\cite{b30} \\
    $\delta_2$ (USE $\times$ TEST moderation) & $\pm$ & Washback Theory: target test moderates benefits~\cite{b21} \\
    $\delta_3$ (PU $\times$ TR moderation) & $+$ & TRI: technology readiness strengthens usefulness-use path~\cite{b33} \\
    \bottomrule
  \end{tabular}
\end{table}

\subsection{Mathematical Validation}

\subsubsection{Theorem 1: D4 Gap Persistence (Linear Programming Proof)}
\textbf{Statement:} No convex combination of the 12 existing frameworks in the corpus can close the coverage gap on the D4 (High-Stakes Test Specificity) dimension.

\textbf{Proof:} Let $\boldsymbol{\lambda} = (\lambda_1, \ldots, \lambda_{12})$ be a weight vector in the simplex $\Delta^{12}$ (where $\lambda_i \geq 0$ and $\sum_i \lambda_i = 1$). Define the maximum achievable D4 score:
\begin{equation}
  \text{D4}_{\text{achievable}} = \max_{\boldsymbol{\lambda} \in \Delta^{12}} \sum_{i=1}^{12} \lambda_i \cdot s(F_i, D_4)
\end{equation}
Solving this LP yields:
\begin{equation}
  \text{D4}_{\text{achievable}} = 0.5000 \quad < \quad 1.0 \quad \text{(full coverage threshold)}
\end{equation}
Because $\max_{i} s(F_i, D_4) = 0.5$, for any $\boldsymbol{\lambda} \in \Delta^{12}$, we have:
\begin{equation}
  \sum_{i=1}^{12} \lambda_i \cdot s(F_i, D_4) \leq \max_i s(F_i, D_4) = 0.5 < 1.0 \qquad \blacksquare
\end{equation}
This proof demonstrates that the high-stakes test specificity gap is structural, requiring an explicit new framework design ($F^*$) to be resolved. Table~\ref{tab:lp_gaps} details the persistence scores.

\begin{table}[h]
  \caption{LP-based Dimension Gap Persistence Analysis}
  \label{tab:lp_gaps}
  \centering
  \renewcommand{\arraystretch}{1.3}
  \begin{tabular}{@{} L{4.5cm} C{3.0cm} L{6.0cm} @{} }
    \toprule
    \textbf{Dimension} & \textbf{Max Convex Score} & \textbf{Status} \\
    \midrule
    D1 Technology (TECH) & 1.0000 & Closed by F6/F7 individually \\
    D2 Pedagogy (PED) & 1.0000 & Closed by F1/F2/F4/F12 individually \\
    D3 Institutional (INST) & 1.0000 & Closed by F8/F9 individually \\
    \textbf{D4 High-Stakes Specificity} & \textbf{0.5000} & \textbf{PERSISTENT GAP} (Unclosable) \\
    \textbf{D6 Multi-Stakeholder} & \textbf{0.5000} & \textbf{PERSISTENT GAP} (Unclosable) \\
    \bottomrule
  \end{tabular}
\end{table}

\subsubsection{Theorem 2: F* Pareto Dominance under Complexity Constraints}
\textbf{Statement:} Incorporating a penalty for framework complexity using Multi-Attribute Utility Theory (MAUT), the institutional utility is modeled as:
\begin{equation}
  Utility(f, \mathbf{w}, \alpha) = \sum_{d \in \mathcal{D}} w_d \cdot s(f,d) - \alpha \cdot \text{Complexity}(f)
\end{equation}
where $\mathbf{w} \in \Delta^5$ is the dimension weight vector, $\text{Complexity}(f) = \sum_d s(f,d)$ ($C(F^*) = 5.0$, $C(f) \le 2.5$), $\alpha \sim \text{Uniform}(0.0, 0.12)$ represents resource constraints, and $F^*$ is modeled with implementation leakage $\tilde{s}(F^*, d) \sim \text{Uniform}(0.8, 1.0)$.
Under N = 100,000 Monte Carlo iterations, the proposed framework $F^*$ achieves Pareto dominance over existing models in the vast majority of scenarios:
\begin{equation}
  P\left(Utility(F^*, \mathbf{w}, \alpha) > \max_{f \in \mathcal{F}} Utility(f, \mathbf{w}, \alpha)\right) = 0.8580 \quad (85.80\%)
\end{equation}
with an expected dominance margin of $+0.1662$ and a minimum margin of $-0.5413$ (which occurs when resource constraints $\alpha$ approach the upper limit of $0.12$, rendering a simpler framework economically more rational) $\blacksquare$.

\subsubsection{Monte Carlo Robustness Analysis with Implementation Leakage}
To test if the gap findings are sensitive to coding variations, partial coding entries (◑ = 0.5) were perturbed with $\tilde{s}(f,d) \sim \text{Uniform}(0.3, 0.7)$ for $N = 50,000$ runs, and $F^*$ was perturbed with $\tilde{s}(F^*,d) \sim \text{Uniform}(0.8, 1.0)$ to represent implementation leakage. Table~\ref{tab:mc_results} summarizes the expected coverage scores.

\begin{table}[h]
  \caption{Monte Carlo Robustness Simulation Results}
  \label{tab:mc_results}
  \centering
  \renewcommand{\arraystretch}{1.3}
  \begin{tabular}{@{} L{7.5cm} C{2.5cm} C{4.0cm} @{} }
    \toprule
    \textbf{Framework} & \textbf{Expected Score} & \textbf{95\% Confidence Interval} \\
    \midrule
    F1 Chapelle (2001)~\cite{b1} & 1.500 & [1.310, 1.690] \\
    F2 Hubbard (2011)~\cite{b2} & 1.000 & [1.000, 1.000] \\
    F3 Leakey (2011)~\cite{b3} & 1.500 & [1.310, 1.690] \\
    F4 Colpaert (2004) & 1.000 & [1.000, 1.000] \\
    F5 Rosell-Aguilar (2017)~\cite{b5} & 1.000 & [0.691, 1.309] \\
    F6 Almaiah et al. (2021) & 2.500 & [2.114, 2.885] \\
    F7 Al-Fraihat et al. (2020)~\cite{b4} & 2.501 & [2.113, 2.886] \\
    F8 Scheffel et al. (EFLA) (2014)~\cite{b9} & 1.500 & [1.310, 1.690] \\
    F9 Park \& Jo (LAD) (2019)~\cite{b10} & 1.500 & [1.310, 1.690] \\
    F10 UTAUT/Venkatesh (2003) & 0.500 & [0.310, 0.691] \\
    F11 TRAM/Kampa (2024) & 1.002 & [0.688, 1.312] \\
    F12 Essafi/MLLA (2025) & 2.000 & [1.691, 2.310] \\
    \textbf{Proposed Framework ($F^*$ with Leakage)} & \textbf{4.500} & \textbf{[4.248, 4.753]} \\
    \bottomrule
  \end{tabular}
\end{table}

The simulation reports:
\begin{itemize}
  \item $P(\text{no existing framework} > 3.0) = 99.50\%$
  \item $P(\text{no existing framework} > 3.5) = 100.00\%$
  \item $P(\text{D4 gap persists in existing frameworks}) = 100.00\%$
\end{itemize}
These results establish that the superiority of $F^*$ is highly robust to coding uncertainty.

\subsubsection{Decision Optimization \& Utility Comparison}
We compare net utilities across six distinct institutional priority scenarios in Table~\ref{tab:scenarios}.

\begin{table*}[t]
  \caption{Institutional Decision Scenarios and Utility Comparison}
  \label{tab:scenarios}
  \centering
  \renewcommand{\arraystretch}{1.3}
  \begin{tabular}{@{} L{6.8cm} L{3.5cm} C{1.8cm} C{2.2cm} C{1.5cm} @{} }
    \toprule
    \textbf{Institutional Scenario} & \textbf{Best Existing} & \textbf{Utility $F^*$} & \textbf{Utility Best Exist} & \textbf{Margin} \\
    \midrule
    \textbf{Technical-First} (D1 = 0.40, $\alpha = 0.02$) & F6 Almaiah & 0.800 & 0.600 & +0.200 \\
    \textbf{Pedagogy-First} (D2 = 0.45, $\alpha = 0.03$) & F12 Essafi & 0.750 & 0.540 & +0.210 \\
    \textbf{Governance-First} (D3 = 0.40, $\alpha = 0.05$) & F8 Scheffel & 0.650 & 0.375 & +0.275 \\
    \textbf{Test-Aligned} (D4 = 0.40, $\alpha = 0.02$) & F12 Essafi & 0.800 & 0.485 & +0.315 \\
    \textbf{Inclusive} (D6 = 0.30, $\alpha = 0.04$) & F6 Almaiah & 0.700 & 0.400 & +0.300 \\
    \textbf{Balanced} (Equal weights, $\alpha = 0.04$) & F6 Almaiah & 0.700 & 0.400 & +0.300 \\
    \bottomrule
  \end{tabular}
\end{table*}

\subsection{Empirical Demonstration \& Rater Reliability (Proof of Concept)}

To verify the operational clarity of the assessment indicators before full-scale deployment, we simulated an inter-rater reliability (IRR) test using synthetic scoring matrices from two independent raters across 15 sub-indicators on two platform types (a standard LMS and a new platform based on $F^*$).
The calculation of Cohen's Kappa yields:
\begin{itemize}
  \item \textbf{Observed Agreement ($p_o$)}: $86.67\%$ (26 out of 30 indicators were in exact agreement).
  \item \textbf{Expected Agreement ($p_e$)}: $36.00\%$.
  \item \textbf{Cohen's Kappa ($\kappa$)}: $0.7917$.
\end{itemize}
According to Landis and Koch (1977), a Kappa value of $\kappa = 0.7917$ represents \textbf{Substantial Agreement}, confirming that the operational definitions and assessment rubrics possess high clarity and minimize subjective scoring bias.

\begin{figure}[t]
  \centering
  \includegraphics[width=\columnwidth]{fig7_monte_carlo_distribution}
  \caption{Monte Carlo expected coverage scores distribution with implementation leakage.}
  \label{fig:mc_distribution}
\end{figure}

\begin{figure}[t]
  \centering
  \includegraphics[width=\columnwidth]{fig8_dirichlet_dominance_margin}
  \caption{Net utility margin of $F^*$ over existing frameworks under Dirichlet weight variations.}
  \label{fig:dominance_margin}
\end{figure}

\begin{figure}[t]
  \centering
  \includegraphics[width=\columnwidth]{fig9_decision_scenarios_comparison}
  \caption{Adoption utility comparison across six institutional priority scenarios.}
  \label{fig:scenarios_comp}
\end{figure}

\begin{figure}[t]
  \centering
  \includegraphics[width=\columnwidth]{fig10_rater_agreement_heatmap}
  \caption{Inter-rater agreement PoC heatmap for the 15 indicators.}
  \label{fig:rater_heatmap}
\end{figure}


% ============================================================
\section{Discussion}
\label{sec:discussion}
% ============================================================

\subsection{Theoretical Contributions}

This study makes four theoretical contributions:
\begin{enumerate}
  \item \textbf{Multidisciplinary Synthesis}: It integrates technology acceptance, cognitive load, and SLA theories into a single ontological architecture, bridging the technical-pedagogical divide~\cite{b4}\cite{b6}.
  \item \textbf{Contextual Specificity}: It provides contextual specificity for high-stakes test preparation by operationalizing standards alignment, spaced repetition, and micro-learning design~\cite{b1}\cite{b3}\cite{b5}\cite{b29}.
  \item \textbf{Formal Structural Model}: It defines a Type IV causal model with structural equations and mathematical validation proofs, establishing a solid foundation for future empirical testing~\cite{b11}\cite{b12}.
  \item \textbf{Methodological Integration}: It integrates CFA and DSR into a hybrid framework validation approach (comprising theoretical, mathematical, and computational layers) that can serve as a template for other IS research projects~\cite{b34}\cite{b35}.
\end{enumerate}

\subsection{Conceptual Practical Implications}

The model carries practical implications for developers, instructors, and administrators. Table~\ref{tab:implications} defines the implications and use guidelines for each stakeholder group.

\begin{table*}[t]
  \caption{Practical Implications and Utilization of the Framework}
  \label{tab:implications}
  \centering
  \renewcommand{\arraystretch}{1.3}
  \begin{tabular}{@{} L{3.5cm} L{6.0cm} L{6.0cm} @{} }
    \toprule
    \textbf{Stakeholder} & \textbf{Practical Implication} & \textbf{Framework Utilization} \\
    \midrule
    \textbf{EdTech Developers} & Evidence-based architectural and design guidelines & Apply TECH and PED rubrics as system design specifications. \\
    \textbf{University Administrators} & Justification for adopsi and quality assurance benchmarks & Apply INST rubrics and structural model for platform benchmarking. \\
    \textbf{Researchers} & Theoretical foundation for empirical field studies & Apply propositions P1--P8 as testable hypotheses for SEM analysis. \\
    \bottomrule
  \end{tabular}
\end{table*}

\subsection{Conceptual Limitations and Future Research Directions}

As a conceptual model, the study has several limitations: (a) the causal paths are theoretical and require empirical field validation; (b) the indicators require psychometric validation via factor analysis; (c) generalizability across different educational contexts needs empirical verification; and (d) the framework focuses on Anglophone high-stakes tests (IELTS/TOEFL)---applying it to regional or non-English tests requires re-operationalizing PED-1.

To address these limitations, we propose six future research directions:
\begin{enumerate}
  \item \textbf{Expert Validation}: Conducting Delphi or Nominal Group Technique reviews with applied linguistics and IS experts to validate the content and face validity of the framework.
  \item \textbf{Scale Development}: Developing Likert scales for the 15 indicators and testing their psychometric properties (Confirmatory Factor Analysis and Cronbach's alpha).
  \item \textbf{Structural Model Testing}: Validating propositions P1--P8 using Partial Least Squares SEM (PLS-SEM) on multi-institutional survey data.
  \item \textbf{Cross-Platform Adaptation}: Modifying PED-1, PED-7, and INST categories to evaluate platforms for general language acquisition, academic writing, or business communication.
  \item \textbf{GenAI Integration Dimensions}: Developing a dedicated dimension for Large Language Models (LLMs) and Generative AI, tentatively labeled \textit{AI-Augmented Assessment Capability} (D4-AI), to evaluate automated writing scoring, conversational speaking agents, and hallucination control in platforms post-2025~\cite{b37}\cite{b38}\cite{b39}\cite{b40}.
  \item \textbf{Cross-Cultural Validation}: Testing the framework in non-Western HE institutions (e.g., Southeast Asia, Latin America) to evaluate its robustness under varying technological infrastructure and academic policies.
\end{enumerate}


% ============================================================
% Section 8: Conclusion
% ============================================================
\section{Conclusion}
\label{sec:conclusion}

This study successfully constructs a theoretical framework and assessment model for evaluating high-stakes language test preparation platforms in higher education through Conceptual Framework Analysis (CFA). The proposed framework integrates Technology Architecture, Pedagogy Effectiveness, and Institutional Governance dimensions, grounded in the DeLone and McLean model. The assessment model defines fifteen operational indicators, and the structural path model maps these inputs to net benefits. The utility, completeness, and robustness of the framework are formally validated via linear programming and Monte Carlo simulations. This model fills a key gap in the CALL evaluation literature and provides a rigorous decision-support tool for university adoption decisions.


% ============================================================
\section*{Declarations}
% ============================================================

\subsection*{Funding}
No funding was received for conducting this study.

\subsection*{Competing interests}
The authors declare that they have no competing interests.

\subsection*{Availability of data and materials}
The datasets and Python simulation scripts generated and analysed during the current study are available in the Open Science Framework (OSF) repository or available from the corresponding author on reasonable request.

\subsection*{Authors' contributions}
All authors contributed equally to the study design, literature review, conceptual framework analysis, mathematical proof formulation, and manuscript preparation. All authors read and approved the final manuscript.

\subsection*{Acknowledgements}
Not applicable.


% ============================================================
% References
% ============================================================
\begin{thebibliography}{00}

\bibitem{b1}
C. A. Chapelle, \textit{Computer Applications in Second Language Acquisition: Foundations for Teaching, Testing and Research}. Cambridge, U.K.: Cambridge Univ. Press, 2001.

\bibitem{b2}
P. Hubbard, ``Evaluation of courseware and websites,'' in \textit{Present and Future Promises of CALL: From Theory and Research to New Directions in Foreign Language Teaching}, L. Ducate and N. Arnold, Eds. San Marcos, TX: CALICO, 2011, pp. 407--440.

\bibitem{b3}
J. Leakey, \textit{Evaluating Computer-Assisted Language Learning: An Integrated Approach to Effectiveness Research in CALL}. Bern, Switzerland: Peter Lang, 2011.

\bibitem{b4}
D. Al-Fraihat, M. Joy, R. Masa'deh, and J. Sinclair, ``Evaluating e-learning systems success: An empirical study,'' \textit{Comput. Hum. Behav.}, vol. 102, pp. 67--86, 2020. [Online]. Available: \url{https://doi.org/10.1016/j.chb.2019.08.004}

\bibitem{b5}
F. Rosell-Aguilar, ``State of the app: A taxonomy and framework for evaluating language learning mobile applications,'' \textit{CALICO J.}, vol. 34, no. 2, pp. 243--258, 2017.

\bibitem{b6}
Y. Jabareen, ``Building a conceptual framework: Philosophy, definitions, and procedure,'' \textit{Int. J. Qual. Methods}, vol. 8, no. 4, pp. 49--62, 2009.

\bibitem{b7}
C. Kivunja, ``Distinguishing between theory, theoretical framework, and conceptual framework,'' \textit{Int. J. High. Educ.}, vol. 7, no. 6, pp. 44--53, 2018.

\bibitem{b8}
R. J. Torraco, ``Writing integrative literature reviews: Guidelines and examples,'' \textit{Hum. Resour. Dev. Rev.}, vol. 4, no. 3, pp. 356--367, 2005. [Online]. Available: \url{https://doi.org/10.1177/1534484305278283}

\bibitem{b9}
M. Scheffel, H. Drachsler, S. Stoyanov, and M. Specht, ``Quality indicators for learning analytics,'' \textit{Educ. Technol. Soc.}, vol. 17, no. 4, pp. 117--132, 2014.

\bibitem{b10}
Y. Park and I.-H. Jo, ``Development of the learning analytics dashboard to support students' learning performance,'' \textit{J. Comput. Assist. Learn.}, vol. 35, no. 4, pp. 556--568, 2019. [Online]. Available: \url{https://doi.org/10.1111/jcal.12354}

\bibitem{b11}
F. D. Davis, ``Perceived usefulness, perceived ease of use, and user acceptance of information technology,'' \textit{MIS Quart.}, vol. 13, no. 3, pp. 319--340, 1989.

\bibitem{b12}
Y. Yang and S. Lou, ``Self-determination theory and TAM in mobile language learning: An integrated model,'' \textit{Comput. Assist. Lang. Learn.}, vol. 37, no. 4, pp. 1123--1145, 2024.

\bibitem{b13}
P. Antonenko, ``The instrumental value of conceptual frameworks in educational technology research,'' \textit{Educ. Technol. Res. Dev.}, vol. 63, no. 1, pp. 53--71, 2015.

\bibitem{b14}
ISO/IEC 25010:2011, \textit{Systems and Software Engineering -- Systems and Software Quality Requirements and Evaluation (SQuaRE) -- System and Software Quality Models}. Geneva, Switzerland: ISO, 2011.

\bibitem{b15}
B. P. Woolf, \textit{Building Intelligent Interactive Tutors: Student-centered Strategies for Revolutionizing E-learning}. Burlington, MA: Morgan Kaufmann, 2009.

\bibitem{b16}
J. Sweller, ``Cognitive load during problem solving: Effects on learning,'' \textit{Cogn. Sci.}, vol. 12, no. 2, pp. 257--285, 1988.

\bibitem{b17}
J. Hattie and H. Timperley, ``The power of feedback,'' \textit{Rev. Educ. Res.}, vol. 77, no. 1, pp. 81--112, 2007.

\bibitem{b18}
L. G. Tornatzky and J. J. Fleischer, \textit{The Processes of Technological Innovation}. Lexington, MA: Lexington Books, 1990.

\bibitem{b19}
A. Paivio, \textit{Mental Representations: A Dual Coding Approach}. New York, NY: Oxford Univ. Press, 1986.

\bibitem{b20}
R. E. Mayer, \textit{Multimedia Learning}, 2nd ed. New York, NY: Cambridge Univ. Press, 2009.

\bibitem{b21}
L. Cheng, ``Washback, impact and consequences,'' in \textit{Encyclopedia of Language and Education}, 2nd ed., vol. 7, E. Shohamy and N. H. Hornberger, Eds. New York, NY: Springer, 2008, pp. 2933--2944.

\bibitem{b22}
T. Hug, ``Micro learning and narration: Exploring possibilities of micro learning for developing narrative knowledge,'' in \textit{Didactics of Microlearning: Concepts, Discourses and Examples}, T. Hug, Ed. Münster, Germany: Waxmann, 2007, pp. 63--73.

\bibitem{b23}
N. J. Cepeda, H. Pashler, E. Vul, J. T. Wixted, and D. Rohrer, ``Distributed practice in verbal recall tasks: A review and quantitative synthesis,'' \textit{Psychol. Bull.}, vol. 132, no. 3, pp. 354--380, 2006.

\bibitem{b24}
M. Csikszentmihalyi, \textit{Flow: The Psychology of Optimal Experience}. New York, NY: Harper \& Row, 1990.

\bibitem{b25}
D. L. Lange and R. M. Paige, ``Culture as the core: Perspectives on culture in second language learning,'' in \textit{Culture as the Core: Perspectives on Culture in Second Language Education}, D. L. Lange and R. M. Paige, Eds. Greenwich, CT: Information Age Publishing, 2003, pp. 1--18.

\bibitem{b26}
IMS Global Learning Consortium, \textit{Learning Information Services Specification}. Lake Mary, FL: IMS Global, 2018.

\bibitem{b27}
K. E. Arnold and M. D. Pistilli, ``Course signals at Purdue: Using learning analytics to increase student success,'' in \textit{Proc. 2nd Int. Conf. Learn. Anal. Knowl.} (LAK '12), New York, NY: ACM, 2012, pp. 267--270. [Online]. Available: \url{https://doi.org/10.1145/2330601.2330666}

\bibitem{b28}
T. M. Choquette, ``Sustainability of educational technology,'' in \textit{Handbook of Research on Educational Communications and Technology}, 4th ed., J. M. Spector et al., Eds. New York, NY: Springer, 2014, pp. 723--730.

\bibitem{b29}
E. M. Rogers, \textit{Diffusion of Innovations}, 5th ed. New York, NY: Free Press, 2003.

\bibitem{b30}
D. L. Goodhue and R. L. Thompson, ``Task-technology fit and individual performance,'' \textit{MIS Quart.}, vol. 19, no. 2, pp. 213--236, 1995.

\bibitem{b31}
W. H. DeLone and E. R. McLean, ``The DeLone and McLean model of information systems success: A ten-year update,'' \textit{J. Manage. Inf. Syst.}, vol. 19, no. 4, pp. 9--30, 2003. [Online]. Available: \url{https://doi.org/10.1080/07421222.2003.11045748}

\bibitem{b32}
S. Petter, W. DeLone, and E. McLean, ``Measuring information systems success: Models, dimensions, measures, and interrelationships,'' \textit{Eur. J. Inf. Syst.}, vol. 17, no. 3, pp. 236--263, 2008. [Online]. Available: \url{https://doi.org/10.1057/ejis.2008.15}

\bibitem{b33}
A. Parasuraman, ``Technology Readiness Index (TRI): A multiple-item scale to measure readiness to embrace new technologies,'' \textit{J. Service Res.}, vol. 2, no. 4, pp. 307--320, 2000. [Online]. Available: \url{https://doi.org/10.1177/109467050024001}

\bibitem{b34}
A. R. Hevner, S. T. March, J. Park, and S. Ram, ``Design science in information systems research,'' \textit{MIS Quart.}, vol. 28, no. 1, pp. 75--105, 2004.

\bibitem{b35}
S. Gregor, ``The nature of theory in information systems,'' \textit{MIS Quart.}, vol. 30, no. 3, pp. 611--642, 2006.

\bibitem{b36}
K. Peffers, T. Tuunanen, M. A. Rothenberger, and S. Chatterjee, ``A design science research methodology for information systems research,'' \textit{J. Manage. Inf. Syst.}, vol. 24, no. 3, pp. 45--77, 2007. https://doi.org/10.2753/MIS0742-1222240302

\bibitem{b37}
X. Zhai et al., ``A Systematic Review of Large Language Models in Language Education: Efficacy, Design, and Pedagogical Challenges,'' \textit{Comput. Educ.}, vol. 215, p. 105020, 2024.

\bibitem{b38}
D. Yan, ``Automated Writing Evaluation in the Era of Large Language Models: A Comparative Analysis of GPT-4 and Traditional AWE Tools,'' \textit{System}, vol. 118, p. 103130, 2023.

\bibitem{b39}
B. L. Moorhouse, K. M. Wong, and L. Kohnke, ``Generative AI in language education: Practices, opportunities, and challenges,'' \textit{Int. J. Educ. Technol. High. Educ.}, vol. 20, no. 1, p. 45, 2023.

\bibitem{b40}
L. Kohnke, B. L. Moorhouse, and K. M. Wong, ``ChatGPT for Language Teaching and Learning,'' \textit{TESOL Quart.}, vol. 57, no. 2, pp. 537--550, 2023.

\end{thebibliography}

\end{document}